\chapter{JUSTIFICATIVA}
\label{capitulo-justificativa}

A mobilidade em ambientes desconhecidos ou com obstáculos pode representar uma dificuldade para pessoas com baixa visão. Nesse contexto, recursos de tecnologia assistiva podem contribuir para aumentar a percepção do ambiente e oferecer informações complementares ao usuário.

A escolha do ESP32 se deve à sua capacidade de processamento, conectividade, baixo custo e possibilidade de integração com sensores e atuadores em um dispositivo compacto. O bracelete proposto pretende detectar situações de proximidade de obstáculos e transformar essas informações em sinais de orientação perceptíveis ao usuário, priorizando uma interação simples e não invasiva.

A realização desta extensão também contribui para a formação acadêmica dos estudantes, permitindo aplicar conhecimentos de programação, eletrônica, sistemas embarcados, sensores e desenvolvimento de protótipos em um problema de relevância social. Caso a ação não seja realizada, perde-se uma oportunidade de experimentar uma solução tecnológica voltada à acessibilidade e de aproximar o conhecimento produzido no curso de uma necessidade da comunidade.